\section{Locomoción robótica: locomoción con patas}

\subsection{Caminata de robots cuadrúpedos}

El ponente presentó casos en los que se usa RL para entrenar robots cuadrúpedos (como Unitree Go1/A1) para caminar sobre distintos terrenos.

\begin{importantbox}{Marco de RL para la locomoción con patas}
\begin{itemize}
    \item \textbf{Estado}: ángulos de las articulaciones, velocidades angulares, datos de la IMU (ángulo de inclinación, velocidad angular)
    \item \textbf{Acción}: ángulo objetivo de cada articulación (seguido por un controlador PD)
    \item \textbf{Recompensa}: velocidad de avance + penalización por energía + penalización por postura + recompensa por frecuencia de pasos
    \item \textbf{Entrenamiento}: se realiza en simuladores paralelos en GPU como Isaac Gym, que pueden ejecutar miles de entornos al mismo tiempo
\end{itemize}
\end{importantbox}

\subsection{Marco Teacher-Student}

\begin{knowledgebox}{Destilación de información privilegiada}
Un paradigma de entrenamiento sim-to-real muy eficaz:
\begin{enumerate}
    \item \textbf{Fase Teacher}: se entrena en simulación una política teacher a la que se le permite acceder a \textbf{información privilegiada} (como el mapa de alturas exacto del terreno, la masa de los objetos y el coeficiente de fricción)
    \item \textbf{Fase Student}: se entrena una política student que solo usa \textbf{sensores disponibles en el mundo real} (como la propiocepción y las cámaras) y que imita el comportamiento del teacher mediante destilación
\end{enumerate}
El teacher puede aprender mejor (porque dispone de más información), mientras que el student aprende a inferir información equivalente a partir de sensores limitados.
\end{knowledgebox}

\begin{figure}[H]
\centering
\includegraphics[width=0.9\textwidth]{slides-images/slide-16.jpg}
\caption{Lo esencial en la fase de despliegue no es «reutilizar directamente la política de simulación», sino estimar en línea los parámetros extrínsecos a partir del historial de observaciones y corregir el controlador de forma continua; esta es también la fuente principal de la robustez de los robots con patas modernos.}
\end{figure}

\subsection{Resumen de la sección}

RL + sim-to-real se ha convertido en el paradigma dominante para el control de robots con patas, y el marco teacher-student resuelve de forma eficaz el problema de la percepción limitada.
